\section{结论与展望}

\subsection{主要发现}

1. \textbf{任务一（月度能耗预测）}：数据探索识别出``能耗 $\propto$ 周里程''的比例结构，单位里程能耗近似与特征无关的随机量。最终以单变量比率均值模型 $energy = 0.81635 \times weekly$ 得到 adj-$R^2=0.876298$。题面建议的``里程 $\times$ 车型''交互与全特征模型均使 adj-$R^2$ 下降。

2. \textbf{任务二（月度充电成本预测）}：充电成本近似由 ``能耗 $\times$ 电价''完全确定（$R^2=0.9994$）。本文以严格嵌套 OOF 级联规避 teacher forcing，最终模型 $cost = 0.81635 \times weekly \times price$ 得到 adj-$R^2=0.917817$，显著优于无级联基线。

3. \textbf{任务三（燃油替代支出预测）}：燃油支出由日通勤里程单变量驱动，叠加标准差约 40 的加性噪声。最终以过原点 OLS $fuel = 8.39476 \times daily$ 得到 adj-$R^2=0.905416$。保留负值样本优于删除、缩尾与 Huber 稳健回归。

4. \textbf{复杂模型全面证伪}：三个任务的浅层 LightGBM 残差模型在 Bootstrap 95\% 置信区间下 $\Delta R^2$ 全部为负，题面提示的交互特征均无增益，验证了``物理结构优先、复杂模型克制''的建模思路。

\subsection{方法学贡献}

\begin{enumerate}[itemsep=0pt]
\item \textbf{数据生成机制识别}：通过比率诊断与相关分析，在建模前即识别出近似确定性的生成公式，为模型选择提供了直接依据；
\item \textbf{严格嵌套级联设计}：任务二的 T1$\to$T2 级联在训练折内部再次交叉验证生成 OOF 能耗，从机制上杜绝 teacher forcing；
\item \textbf{证据驱动的模型晋级门槛}：以 pooled OOF adj-$R^2$ 增益、Bootstrap 95\% 置信区间与多随机种子稳定性三重标准判定模型取舍，而非依赖单次实验的偶然提升；
\item \textbf{负结果的诚实报告}：完整记录交互特征与复杂模型被证伪的全过程，为同类问题提供可复用的信号验证模板。
\end{enumerate}

\subsection{局限性}

1. \textbf{口径依赖组委会确认}：adj-$R^2$ 中自由度 $p$ 的定义、正式测试集的字段是否隐藏三个目标列、T2 是否提供真实能耗等问题尚待组委会明确，本文采用最保守的严格口径。

2. \textbf{非线性模型覆盖有限}：本文仅以浅层 LightGBM 作为非线性代表进行残差诊断，未穷举全部可能模型。但从消融与 Bootstrap 结果看，数据结构本身不支持非线性增益，该局限不影响结论。

3. \textbf{公式外推性未知}：最终模型为闭式公式，其系数依赖当前训练集；若正式测试集的生成机制与训练集不同（例如新增了与能效相关的变量），模型的鲁棒性有待检验。

4. \textbf{业务可解释性的边界}：任务三的负值燃油支出样本虽在统计上应保留，但在业务上仍属异常，其处理方式需在报告中向评审说明。
